% !TeX root = ../main.tex
\begin{abstract}
Music-to-text choreography represents a dance as an inspectable, timed
movement score. It requires actionable local descriptions, coherent
whole-song organization, and evaluation that admits multiple valid dances
for the same music. We introduce \method through an evidence-linked
annotation resource, a dedicated model trained by supervised fine-tuning,
a compositional harness, and a multi-level evaluation protocol. The
design decouples local movement prediction from global composition:
short-clip proposals are organized into bar-aligned units using musical
evidence, global style, segment objectives, and motif memory. Schedule
validation and bounded revision preserve accepted context. Source-event
links make annotation refinement auditable. On a frozen 2,000-clip panel, the adapted
model achieves 30.80 ROUGE-L versus 14.99 for the best evaluated general
audio-language baseline on this metric, demonstrating stronger reference
agreement. Our evaluation protocol separately addresses local musical
response and whole-song organization. Preliminary whole-song results
concern the historical fixed-window harness; controlled comparison of
the latest bar-aligned agents and independent audio-grounded assessment
remain prospective.
\end{abstract}
